
\begin{proposicion}[Un menor lejano; estabilidad \textup{(S)}]\label{prop:menor-lejano}
Para constantes $(C_p,\gamma_A)$ hay un umbral $k_S=k_S(C_p,\gamma_A)\ge k_2$ y una constante
$C_{S1}=C_{S1}(C_p,\gamma_A)$ con las propiedades siguientes.
\begin{enumerate}
\item[\textup{(a)}] Sea $N$ una máscara que cumple \textup{(H)} con $(C_p,\gamma_A)$, de tamaño
$k\ge k_S$, sea $d$ una casilla de $N$ y $c$ una casilla de $N_d$. Entonces
$|r_c(N_d)-r_c(N)|\le C_{S1}/k^3$.
\item[\textup{(b)}] Sea $M$ una máscara que cumple \textup{(H)} con $(C_p,\gamma_A)$, de tamaño $k$, con
$k\ge k_S(C_p,\gamma_A)$ y $k-1\ge k_S(\tfrac54C_p,\tfrac34\gamma_A)$. Sean $e,f$ casillas en filas y
columnas distintas y $c$ una casilla que no está en ninguna de sus filas ni de sus columnas. Entonces
\[
 |r_c(M_{ef})-r_c(M)|\le\frac{C_S}{k^3}, \tag{S}
\]
con $C_S$ que sólo depende de $(C_p,\gamma_A)$.
\end{enumerate}
En la clase $\mathcal K$ ambas valen con $(C_p,\gamma_A)=(\tfrac32M_0,0{,}88\gamma)$.
\end{proposicion}
\begin{proof}
La parte (b) se sigue de (a). En efecto, $M_{ef}=(M_e)_f$; por el Paso~4 de la demostración del
Teorema~\ref{teo:insercion-torres}, $M_e$ cumple \textup{(H)} con $(\tfrac54C_p,\tfrac34\gamma_A)$ en
tamaño $k-1$; y $c$ es una casilla de $M_e$ y de $M_{ef}$. Se aplica (a) a $(M,e)$ y a $(M_e,f)$ y se
suma.

Para (a) usamos la notación de aquella demostración con $N$ en lugar de $M$ y $e_0=c$, escribimos
$h=1/k$, y $\delta$ denota la variación de una cantidad al pasar de $N$ a $N_d$, leída en el soporte
común. Las constantes sólo dependen de $(C_p,\gamma_A)$, y $k_S$ se elige de modo que valgan las
desigualdades elementales que se usan abajo.

\emph{Paso 1. Un representante de segundo orden de $r_{e_0}$.} Sean $u_x=a_x^\top\xi_N$,
$\alpha_x=\mathsf G(x,x)-2$, $\eta_x=\mathsf G(x,e_0)$, $b_x=\tfrac12\tau_x\dot\tau_xu_x^2$ con
$\dot\tau_x=\dot\tau(p_x)$, y
\[
 Z=(S'(\tau|)+\Pi')^{-1}\Bigl(\rho-\sum_{x\in M'}b_xa_x\Bigr),\qquad z_x=a_x^\top Z .
\]
Todo ello se calcula con los datos del padre. La ecuación de masas del hijo es exactamente
\begin{equation}\label{eq:masas-hijo}
 \zeta=(S'(\tau|)+\Pi')^{-1}(\rho-B),\qquad
 B_i=\sum_{x\in i\cap M'}\bigl[u(t_xe^{s_x})-p_x-\tau_xs_x\bigr].
\end{equation}
Como $\|\xi\|_\infty=O(h)$ y $\|\xi-\xi_N\|_\infty=O(h^2)$ por~\eqref{eq:newton-hijo}, la tercera
derivada de $s\mapsto u(te^s)$ es $O(p)$ en el segmento, las masas de línea están acotadas y la inversa
está acotada en la norma del supremo, un desarrollo de Taylor da $\zeta=Z+O_\infty(h^3)$.

En~\eqref{eq:re-restos} reemplazamos los restos como sigue; $E_{\mathcal D3}=\zeta^\top\rho$ y
$E_{\mathcal D4}=-\tfrac12\zeta^\top S'(\tau|)\zeta$ ya son $O(h^3)$.
\begin{enumerate}
\item En $E_{\rm main}$ se reemplaza $p_x^3/(1-2p_x)$ por $p_x^3$.
\item Se reemplaza $E_{\mathcal D1}=\tfrac12\sum_{F^*}\varsigma_x\beta_{\lambda(x)}$ por
$\tfrac12\sum_{F^*}p_x^2\beta_{\lambda(x)}$.
\item En $E_{\mathcal D2}=-\tfrac12\beta_I^\top\rho$ se reemplaza $\rho_{\lambda(x)}$ por
$-p_x^2-p_x(\eta_x-1)$ para $x\in F^*$.
\item Se reemplaza $E_{\mathcal D5}=-R_3$ por $-\tfrac16\sum_{M'}\tau_x\dot\tau_xu_x^3$.
\item Se reemplaza $E_{L1}$ por $-\tfrac23\sum_{F^*}p_x^3-p_0^2-\tfrac14(\mu_r^2+\mu_c^2)$.
\item Se reemplaza $E_{L2}$ por $\sum_{M'}\bigl\{\tfrac12D(\ell_2'\tau)_xu_x^2
+(\ell_2'(p_x)-p_x)\tau_xu_x+p_x\tau_xz_x\bigr\}$, donde $\ell_2'(u)=u/[(1-2u)(1-u)]$ y
$D=t\,d/dt$, de modo que $Du=\tau$.
\item Se reemplaza $E_{\rm(i)}$ por $-\sum_{M'}\bigl\{\tau\alpha z+4p\tau z+b\alpha
+2(\dot\varsigma-2p)\tau u+2\dot\varsigma b+\ddot\varsigma(\tau u)^2+\dot\varsigma\tau u\alpha\bigr\}_x$.
\item Se reemplaza $E_{L3}$ por $\tfrac14\operatorname{tr}(\mathsf S^+W_0\mathsf S^+W_0)$, donde $W_0$
se obtiene de las entradas de $W$ del Paso~7 poniendo $\Delta\varsigma_x=\dot\varsigma_x\tau_xu_x$ en
la diagonal y $\tau'_x-\tau_x=\dot\tau_x\tau_xu_x$ fuera de ella.
\item Se reemplaza $E_{\delta1}$ por $\tfrac12\hat c_{e_0}^2$.
\item Se reemplaza $E_{\delta2}=p_0\hat c_{e_0}/(1-p_0)$ por $p_0\hat c_{e_0}$.
\item En
$E_X=\sum_{M'}\tau_x\tilde\varepsilon_x\eta_x+\sum_{F^*}\tau_x\tilde\varepsilon_x
+\sum_{F^*}\tau_x\varepsilon_x(\eta_x-1)+\tau_0\varepsilon_0\mathsf G_0-\tilde\varepsilon_0$ se
reemplazan $\tilde\varepsilon_x$ por $6p_x^2+2p_x\alpha_x$ y $\varepsilon_x$ por $\alpha_x+2p_x$.
\end{enumerate}
Sea $r^{[2]}_{e_0}$ la suma de estas expresiones con los coeficientes de~\eqref{eq:re-restos}. Cada
reemplazo cuesta $O(h^3)$, de modo que $r_{e_0}=r^{[2]}_{e_0}+O(h^3)$ uniformemente, en $N$ y en $N_d$.
Para verlo, nótese que en $F^*$ se tiene $\sum p\le2$ y $\sum p^j=O(h^{j-1})$, y en $M'$ se tiene
$\sum p\le k$, $\sum p^2=O(1)$, $\sum p^3=O(h)$. Los errores de los puntos 1 a 3 son
$O(\sum_{F^*}p^4)$ y $O(h\sum_{F^*}(p^3+hp^2))$, y el del punto~4 es $O(h^4\sum_{M'}p)$. Para el
punto~5, $-\log(1-u)-u+\tfrac12c_2(u)=-\tfrac23u^3+O(u^4)$; los términos $\tfrac12c_2(p_0)$ y
$\tfrac12\log(1-\tau_0^2/(m_rm_c))$ aportan $-\tfrac12p_0^2+O(h^3)$ cada uno, y los logaritmos de las dos
masas dan $-\tfrac14(\mu_r^2+\mu_c^2)+O(h^3)$. Para el punto~6, $D(\ell_2'\tau)$ y $D^2(\ell_2'\tau)$ son
$O(p^2)$ y $(\ell_2'-p)\tau=O(p^3)$. Para el punto~7, con $s=a_x^\top\xi$,
\[
 \Delta p-\tau s-b=O(ph^3),\qquad \Delta\varsigma-\dot\varsigma\tau u=O(p^2h^2),\qquad
 \Delta\varsigma-\dot\varsigma\Delta p-\tfrac12\ddot\varsigma(\tau u)^2=O(p^2h^3),
\]
y se suma usando $\dot\varsigma=O(p)$, $\dot\varsigma-2p=O(p^2)$, $\alpha=O(h)$ y $|s-u|=O(h^2)$. Para
el punto~8 el error en la diagonal es $O(h^2\sum_{x\in i}p_x^2)=O(h^3)$ por línea y $O(ph^2)$ por
entrada fuera de la diagonal, luego $\|W-W_0\|_{\rm HS}=O(h^2)$; como $\|W\|_{\rm HS}=O(h)$, la traza
cuadrática cambia en $O(h^3)$, y la cola del logaritmo es $O(h^3)$ por la cota espectral. Los puntos 9
y 10 se siguen de $-\log(1+x)+x=\tfrac12x^2+O(x^3)$ y de $\hat c_{e_0}=O(h)$. Para el punto~11,
$\dot c_2+\dot\tau(2+\alpha)-2=\alpha+2p+6p^2+2p\alpha+O(p^3+p^2|\alpha|)$, y las sumas se controlan con
$\sum_{M'}\tau|\eta|\le C$, $|\eta-1|\le Ch$ en $F^*$ y $\tau_0|\mathsf G_0|=O(h)$.

\emph{Paso 2. Variación del desplazamiento.} El Teorema~\ref{teo:insercion-torres} y el
Corolario~\ref{cor:dos-casillas}, aplicados al menor de $d$, dan en el soporte común
\begin{equation}\label{eq:variaciones-base}
 \delta p=O(ph),\qquad \delta G^0_{ij}=O(h^2),\qquad \delta\hat c=O(h^2),\qquad
 \delta\tau=O(ph),\qquad \delta\mu_i=O(h^2),
\end{equation}
la última por suma en una línea que pierde a lo sumo un sitio. De $\zeta=O(h^2)$ y
$s_x=-\mathsf G(x,e_0)+a_x^\top\zeta$ resulta primero $\delta s=O(h^2)$. En~\eqref{eq:masas-hijo} el
sumando de $B_i$ es $p_xs_x^2$ por una función con derivadas acotadas en $(p_x,s_x)$ cerca del origen.
En los índices comunes cambia en $O(p_xh^3)$, y el sitio perdido aporta $O(h^3)$; así $\delta
B_i=O(h^3)$. En la estrella $F^*$, la fórmula exacta~\eqref{eq:rho-hijo} da $\delta\rho_i=O(h^3)$. De los
dos sistemas lineales~\eqref{eq:masas-hijo}, para $N$ y para $N_d$, se restringe el mayor a los índices
comunes. Las sumas absolutas por fila de la diferencia de las dos matrices, incluido el acoplamiento
con los dos índices retirados, son $O(h)$: esto se sigue de $\delta\tau=O(ph)$, del sitio perdido y del
cambio de la proyección del calibre. Multiplicando por la solución, que es $O(h^2)$, y por la inversa,
uniformemente acotada en la norma del supremo, obtenemos
\[
 \delta\zeta=O(h^3),\qquad\text{y del mismo modo}\qquad \delta z=O(h^3),\qquad \delta u=O(h^2).
\]

\emph{Paso 3. Variación del resto.} Por el Paso~1 basta comparar las once expresiones en $N$ y en
$N_d$. Los términos interiores de los puntos 4, 6 y 7 tienen mayorantes $Cph^3$, $Cp^2h^2$ y $Cp^3h$
por sitio; cada uno suma $O(h^2)$, y $O(h^3)$ en las dos líneas perdidas. En los índices comunes el
telescopio de cada producto gana un factor $h$ en el mayorante, por~\eqref{eq:variaciones-base},
$\delta\alpha=O(h^2)$, $\delta z=O(h^3)$ y $\delta u=O(h^2)$; los factores escalares de orden $p^m$,
$m\ge1$, son analíticos con derivada $O(p^{m-1})$. En la estrella se pierden a lo sumo dos términos de
tamaño $O(h^3)$. Los productos con $\mu$ usan $\mu=O(h)$ y $\delta\mu=O(h^2)$; los del sitio elegido
usan $\hat c=O(h)$ y $\delta\hat c=O(h^2)$. En el punto~11, $\delta(6p^2+2p\alpha)=O(h^3)$; la suma
interior usa $\sum\tau|\mathsf G(x,e_0)|=O(1)$, y en los sitios retirados la masa es a lo sumo dos y
$\mathsf G(x,e_0)=O(h)$.

Para la traza se puede mantener $W$ exacta, con un coste $O(h^3)$ en cada tablero. Como $W$ anula el
calibre, la traza cuadrática es $\operatorname{tr}(G^0_{II}WG^0_{II}W)$. Las entradas de $W$ del
Paso~7 dan $\max|W|=O(h^2)$, $\|W\|_\infty=O(h)$ y, en los índices comunes, $\max|\delta W|=O(h^3)$ y
$\|\delta W\|_\infty=O(h^2)$; aquí se usan $\delta s=O(h^2)$, $\delta p=O(ph)$ y la fórmula de $Q_J$.
Escribamos $G^0_{II}=I+H_1$, con $\max|H_1|=O(h)$ y $\max|\delta H_1|=O(h^2)$. En
$\operatorname{tr}(W^2)$, cambiar $W$ cuesta $O(k\,h^2h^2)$ y perder una línea cuesta $O(h\,h^2)$. La
suma absoluta de $\operatorname{tr}(H_1W^2)$ es a lo sumo $k\max|H_1|\,\|W\|_\infty^2$, y la de
$\operatorname{tr}(H_1WH_1W)$ es a lo sumo $k^2\max|H_1|^2\|W\|_\infty^2$; cambiar un factor o fijar un
índice perdido mejora cada mayorante en un factor $h$. En total $\delta r_{e_0}=O(h^3)$, que es (a).
\end{proof}

\begin{corolario}[Incremento del resto de dos casillas]\label{cor:incremento-marginal}
En las hipótesis de la Proposición~\ref{prop:menor-lejano}\textup{(b)}, sea $R_\pi=\pi^T-\pi_1$ para el
par $\{e,f\}$. Si $c$ no está en ninguna de las filas ni de las columnas de $e,f$, entonces
\[
 |R_\pi(M_c)-R_\pi(M)|\le C\,\frac{p_ep_f}{k^3}.
\]
\end{corolario}
\begin{proof}
Por~\eqref{eq:psi-dos-casillas}, $R_\pi=p_ep_f\Psi$ con
$\Psi=(1+x_1)(1-g+t_1)(1+x_2+t_2)e^{\varrho_1+\varrho_2}-(1-g+x_1+x_2)$, donde $x_1=\hat c_e$,
$x_2=\hat c_f$, $g=\mathsf G(e,f)$, $t_1=p'_f/p_f-1+g$, $t_2=\hat c'_f-\hat c_f$, $\varrho_1=r_e(M)$ y
$\varrho_2=r_f(M_e)$. Aquí $x_i,g=O(h)$ y $t_i,\varrho_i=O(h^2)$. La parte de $\Psi$ lineal en
$(x_1,x_2,g)$ se anula, luego sus derivadas en $x_i,g$ son $O(h)$, mientras que las derivadas en
$t_i,\varrho_i$ son $1+O(h)$. Sea $\delta$ la variación de $M$ a $M_c$. Entonces
$\delta x_i,\delta g=O(h^2)$ por~\eqref{eq:variaciones-base}, y
$\delta\varrho_1,\delta\varrho_2=O(h^3)$ por la Proposición~\ref{prop:menor-lejano}\textup{(a)},
aplicada a $M$ y a $M_e$. Para $t_1,t_2$ hacen falta segundas diferencias. La matriz de Green del hijo
cumple $G'-G^0_{II}=(I+H_1)W(I+H_2)+C_q$, donde las entradas de $H_i$ son $O(h)$ con variaciones
$O(h^2)$, y $C_q$ se construye con $q=G^0_{IJ}v_J=O(h)$, $\delta q=O(h^2)$, $v_I^\top q=O(h)$ y el
divisor $2k-2$, de modo que $\delta C_q=O(h^3)$. Con las cotas de $W$ y $\delta W$ del Paso~3 resulta
$\delta(G'-G^0_{II})_{ij}=O(h^3)$. Para $X_x=(G^0a_x)^\top\eta^X-\varepsilon_x$ con
$\eta^X_i=\sum_{y\in i}\tau_y\varepsilon_y$, se tiene $\delta\Delta p=O(ph^2)$ por
$\Delta p=\tau s+O(ps^2)$ y~\eqref{eq:masas-hijo}, luego $\delta\Delta\varepsilon=O(h^3)$ y
$\delta\Delta\eta^X_i=O(h^3)$; en la contracción con $G^0a_x$, cuya norma $\ell^1$ es $O(1)$ y cuyas
entradas en las cuatro líneas retiradas son $O(h)$, esto da $\delta\Delta X=O(h^3)$ y
$\delta t_2=O(h^3)$. Por último,
$t_1=-(\varsigma_f/p_f)\mathsf G(e,f)+(\tau_f/p_f)a_f^\top\zeta+[u(t_fe^{s_f})-p_f-\tau_fs_f]/p_f$ tiene
variación $O(h^3)$ por el Paso~2. Por tanto $\delta\Psi=O(h^3)$. Como $\Psi=O(h^2)$ y
$\delta(p_ep_f)=O(p_ep_fh)$, se obtiene la tesis.
\end{proof}

